\begin{frame}{그림으로: RAG 원 논문의 구조}
  \begin{center}
    \includegraphics[width=0.8\linewidth]{rag_fig1_lewis2020.png}
  \end{center}
  {\footnotesize Lewis et al. (2020), \textit{Retrieval-Augmented
  Generation for Knowledge-Intensive NLP Tasks} --- Query Encoder로
  질문을 벡터 $q(x)$로, Document Index에서 코사인/내적(MIPS)으로 가장
  가까운 문서 $z_1,\ldots,z_k$를 찾고(\kb{Retriever}), 이를 생성 모델
  $p_\theta$(\kb{Generator})에 넘겨 답을 만든다 --- 원 논문 Figure 1.}
\end{frame}
